% PhysiGym framework dataflow figure.
% % PhysiGym framework dataflow figure.
% % PhysiGym framework dataflow figure.
% % PhysiGym framework dataflow figure.
% \input{fig_framework} inside a figure environment, OR use as-is (it wraps itself).
% Requires: \usetikzlibrary{arrows.meta, positioning, fit, backgrounds}
\begin{figure}[t]
  \centering
  \begin{tikzpicture}[
      font=\small,
      box/.style={draw, rounded corners, minimum height=8mm, inner sep=4pt, align=center},
      cpp/.style={box, fill=blue!8},
      py/.style={box, fill=green!8},
      gym/.style={box, fill=orange!8},
      flow/.style={-{Latex[length=2mm]}, thick},
      readback/.style={-{Latex[length=2mm]}, thick, dashed},
    ]
    % --- C++ / PhysiCell layer ---
    \node[cpp] (start)  {\texttt{physicell\_start}};
    \node[cpp, right=8mm of start] (step) {\texttt{physicell\_step}\\{\scriptsize(inject \texttt{drug\_1})}};
    \node[cpp, right=8mm of step]  (stop) {\texttt{physicell\_stop}};
    \begin{pgfonlayer}{background}
      \node[draw, rounded corners, fill=blue!8, fit=(start)(step)(stop),
            inner sep=6mm,
            label=above:{\textbf{\physicell{} core (C++) --- BioFVM + ABM}}] (cppbox) {};
    \end{pgfonlayer}

    % --- read functions ---
    \node[cpp, below=10mm of step, xshift=-16mm] (getcell)    {\texttt{get\_cell}};
    \node[cpp, right=4mm of getcell] (getmicro) {\texttt{get\_microenv}};
    \node[cpp, right=4mm of getmicro] (getgraph) {\texttt{get\_graph}};

    % --- Python env layer ---
    \node[py, below=12mm of getmicro] (env) {\texttt{ModelPhysiCellEnv}\\{\scriptsize subclass of \texttt{CorePhysiCellEnv}}\\{\scriptsize get\_observation / get\_reward / get\_action\_space}};

    % --- Gymnasium / agent layer ---
    \node[gym, below=10mm of env] (agent) {SAC agent (Gymnasium API)\\{\scriptsize \texttt{env.step(a)} $\to$ \texttt{(obs, reward, done, info)}}};

    % flows
    \draw[flow] (getcell) -- (env);
    \draw[flow] (getmicro) -- (env);
    \draw[flow] (getgraph) -- (env);
    \draw[readback] (step.south) -- (getmicro.north);
    \draw[flow] (env) -- node[right]{\scriptsize obs, reward} (agent);
    \draw[flow] (agent.west) to[out=180,in=180] node[left,align=right]{\scriptsize action\\\scriptsize (dose,$x$,$y$,$r$)} (step.west);
  \end{tikzpicture}
  \caption{\physigym{} dataflow. The \physicell{} C++ loop is exposed to Python as
    three callables (\texttt{physicell\_start/step/stop}); only \texttt{physicell\_step}
    is modified per model (here, to inject \texttt{drug\_1}). Three read functions
    (\texttt{get\_cell}, \texttt{get\_microenv}, \texttt{get\_graph}) surface cell
    positions, substrate fields, and neighbourhood graphs. A user subclasses
    \texttt{CorePhysiCellEnv} to assemble these into a Gymnasium observation/reward,
    and any RL algorithm drives the environment through the standard
    \texttt{step}/\texttt{reset} API.}
  \label{fig:framework}
\end{figure}
 inside a figure environment, OR use as-is (it wraps itself).
% Requires: \usetikzlibrary{arrows.meta, positioning, fit, backgrounds}
\begin{figure}[t]
  \centering
  \begin{tikzpicture}[
      font=\small,
      box/.style={draw, rounded corners, minimum height=8mm, inner sep=4pt, align=center},
      cpp/.style={box, fill=blue!8},
      py/.style={box, fill=green!8},
      gym/.style={box, fill=orange!8},
      flow/.style={-{Latex[length=2mm]}, thick},
      readback/.style={-{Latex[length=2mm]}, thick, dashed},
    ]
    % --- C++ / PhysiCell layer ---
    \node[cpp] (start)  {\texttt{physicell\_start}};
    \node[cpp, right=8mm of start] (step) {\texttt{physicell\_step}\\{\scriptsize(inject \texttt{drug\_1})}};
    \node[cpp, right=8mm of step]  (stop) {\texttt{physicell\_stop}};
    \begin{pgfonlayer}{background}
      \node[draw, rounded corners, fill=blue!8, fit=(start)(step)(stop),
            inner sep=6mm,
            label=above:{\textbf{\physicell{} core (C++) --- BioFVM + ABM}}] (cppbox) {};
    \end{pgfonlayer}

    % --- read functions ---
    \node[cpp, below=10mm of step, xshift=-16mm] (getcell)    {\texttt{get\_cell}};
    \node[cpp, right=4mm of getcell] (getmicro) {\texttt{get\_microenv}};
    \node[cpp, right=4mm of getmicro] (getgraph) {\texttt{get\_graph}};

    % --- Python env layer ---
    \node[py, below=12mm of getmicro] (env) {\texttt{ModelPhysiCellEnv}\\{\scriptsize subclass of \texttt{CorePhysiCellEnv}}\\{\scriptsize get\_observation / get\_reward / get\_action\_space}};

    % --- Gymnasium / agent layer ---
    \node[gym, below=10mm of env] (agent) {SAC agent (Gymnasium API)\\{\scriptsize \texttt{env.step(a)} $\to$ \texttt{(obs, reward, done, info)}}};

    % flows
    \draw[flow] (getcell) -- (env);
    \draw[flow] (getmicro) -- (env);
    \draw[flow] (getgraph) -- (env);
    \draw[readback] (step.south) -- (getmicro.north);
    \draw[flow] (env) -- node[right]{\scriptsize obs, reward} (agent);
    \draw[flow] (agent.west) to[out=180,in=180] node[left,align=right]{\scriptsize action\\\scriptsize (dose,$x$,$y$,$r$)} (step.west);
  \end{tikzpicture}
  \caption{\physigym{} dataflow. The \physicell{} C++ loop is exposed to Python as
    three callables (\texttt{physicell\_start/step/stop}); only \texttt{physicell\_step}
    is modified per model (here, to inject \texttt{drug\_1}). Three read functions
    (\texttt{get\_cell}, \texttt{get\_microenv}, \texttt{get\_graph}) surface cell
    positions, substrate fields, and neighbourhood graphs. A user subclasses
    \texttt{CorePhysiCellEnv} to assemble these into a Gymnasium observation/reward,
    and any RL algorithm drives the environment through the standard
    \texttt{step}/\texttt{reset} API.}
  \label{fig:framework}
\end{figure}
 inside a figure environment, OR use as-is (it wraps itself).
% Requires: \usetikzlibrary{arrows.meta, positioning, fit, backgrounds}
\begin{figure}[t]
  \centering
  \begin{tikzpicture}[
      font=\small,
      box/.style={draw, rounded corners, minimum height=8mm, inner sep=4pt, align=center},
      cpp/.style={box, fill=blue!8},
      py/.style={box, fill=green!8},
      gym/.style={box, fill=orange!8},
      flow/.style={-{Latex[length=2mm]}, thick},
      readback/.style={-{Latex[length=2mm]}, thick, dashed},
    ]
    % --- C++ / PhysiCell layer ---
    \node[cpp] (start)  {\texttt{physicell\_start}};
    \node[cpp, right=8mm of start] (step) {\texttt{physicell\_step}\\{\scriptsize(inject \texttt{drug\_1})}};
    \node[cpp, right=8mm of step]  (stop) {\texttt{physicell\_stop}};
    \begin{pgfonlayer}{background}
      \node[draw, rounded corners, fill=blue!8, fit=(start)(step)(stop),
            inner sep=6mm,
            label=above:{\textbf{\physicell{} core (C++) --- BioFVM + ABM}}] (cppbox) {};
    \end{pgfonlayer}

    % --- read functions ---
    \node[cpp, below=10mm of step, xshift=-16mm] (getcell)    {\texttt{get\_cell}};
    \node[cpp, right=4mm of getcell] (getmicro) {\texttt{get\_microenv}};
    \node[cpp, right=4mm of getmicro] (getgraph) {\texttt{get\_graph}};

    % --- Python env layer ---
    \node[py, below=12mm of getmicro] (env) {\texttt{ModelPhysiCellEnv}\\{\scriptsize subclass of \texttt{CorePhysiCellEnv}}\\{\scriptsize get\_observation / get\_reward / get\_action\_space}};

    % --- Gymnasium / agent layer ---
    \node[gym, below=10mm of env] (agent) {SAC agent (Gymnasium API)\\{\scriptsize \texttt{env.step(a)} $\to$ \texttt{(obs, reward, done, info)}}};

    % flows
    \draw[flow] (getcell) -- (env);
    \draw[flow] (getmicro) -- (env);
    \draw[flow] (getgraph) -- (env);
    \draw[readback] (step.south) -- (getmicro.north);
    \draw[flow] (env) -- node[right]{\scriptsize obs, reward} (agent);
    \draw[flow] (agent.west) to[out=180,in=180] node[left,align=right]{\scriptsize action\\\scriptsize (dose,$x$,$y$,$r$)} (step.west);
  \end{tikzpicture}
  \caption{\physigym{} dataflow. The \physicell{} C++ loop is exposed to Python as
    three callables (\texttt{physicell\_start/step/stop}); only \texttt{physicell\_step}
    is modified per model (here, to inject \texttt{drug\_1}). Three read functions
    (\texttt{get\_cell}, \texttt{get\_microenv}, \texttt{get\_graph}) surface cell
    positions, substrate fields, and neighbourhood graphs. A user subclasses
    \texttt{CorePhysiCellEnv} to assemble these into a Gymnasium observation/reward,
    and any RL algorithm drives the environment through the standard
    \texttt{step}/\texttt{reset} API.}
  \label{fig:framework}
\end{figure}
 inside a figure environment, OR use as-is (it wraps itself).
% Requires: \usetikzlibrary{arrows.meta, positioning, fit, backgrounds}
\begin{figure}[t]
  \centering
  \begin{tikzpicture}[
      font=\small,
      box/.style={draw, rounded corners, minimum height=8mm, inner sep=4pt, align=center},
      cpp/.style={box, fill=blue!8},
      py/.style={box, fill=green!8},
      gym/.style={box, fill=orange!8},
      flow/.style={-{Latex[length=2mm]}, thick},
      readback/.style={-{Latex[length=2mm]}, thick, dashed},
    ]
    % --- C++ / PhysiCell layer ---
    \node[cpp] (start)  {\texttt{physicell\_start}};
    \node[cpp, right=8mm of start] (step) {\texttt{physicell\_step}\\{\scriptsize(inject \texttt{drug\_1})}};
    \node[cpp, right=8mm of step]  (stop) {\texttt{physicell\_stop}};
    \begin{pgfonlayer}{background}
      \node[draw, rounded corners, fill=blue!8, fit=(start)(step)(stop),
            inner sep=6mm,
            label=above:{\textbf{\physicell{} core (C++) --- BioFVM + ABM}}] (cppbox) {};
    \end{pgfonlayer}

    % --- read functions ---
    \node[cpp, below=10mm of step, xshift=-16mm] (getcell)    {\texttt{get\_cell}};
    \node[cpp, right=4mm of getcell] (getmicro) {\texttt{get\_microenv}};
    \node[cpp, right=4mm of getmicro] (getgraph) {\texttt{get\_graph}};

    % --- Python env layer ---
    \node[py, below=12mm of getmicro] (env) {\texttt{ModelPhysiCellEnv}\\{\scriptsize subclass of \texttt{CorePhysiCellEnv}}\\{\scriptsize get\_observation / get\_reward / get\_action\_space}};

    % --- Gymnasium / agent layer ---
    \node[gym, below=10mm of env] (agent) {SAC agent (Gymnasium API)\\{\scriptsize \texttt{env.step(a)} $\to$ \texttt{(obs, reward, done, info)}}};

    % flows
    \draw[flow] (getcell) -- (env);
    \draw[flow] (getmicro) -- (env);
    \draw[flow] (getgraph) -- (env);
    \draw[readback] (step.south) -- (getmicro.north);
    \draw[flow] (env) -- node[right]{\scriptsize obs, reward} (agent);
    \draw[flow] (agent.west) to[out=180,in=180] node[left,align=right]{\scriptsize action\\\scriptsize (dose,$x$,$y$,$r$)} (step.west);
  \end{tikzpicture}
  \caption{\physigym{} dataflow. The \physicell{} C++ loop is exposed to Python as
    three callables (\texttt{physicell\_start/step/stop}); only \texttt{physicell\_step}
    is modified per model (here, to inject \texttt{drug\_1}). Three read functions
    (\texttt{get\_cell}, \texttt{get\_microenv}, \texttt{get\_graph}) surface cell
    positions, substrate fields, and neighbourhood graphs. A user subclasses
    \texttt{CorePhysiCellEnv} to assemble these into a Gymnasium observation/reward,
    and any RL algorithm drives the environment through the standard
    \texttt{step}/\texttt{reset} API.}
  \label{fig:framework}
\end{figure}
